\documentclass[11pt]{article}
\usepackage[margin=2.6cm]{geometry}
\usepackage{amsmath,amssymb,amsthm,mathtools}
\usepackage{enumitem}
\usepackage{booktabs}
\usepackage{array}
\usepackage{longtable}
\usepackage[colorlinks=true,allcolors=blue]{hyperref}
\usepackage[capitalize,noabbrev]{cleveref}
\usepackage{xspace}

% ------------------------- theorem environments -------------------------
\theoremstyle{remark}
\newtheorem{observation}{Observation}
\newtheorem*{note}{Note}

% ------------------------- macros (mirroring osgs_convergence.tex) -------------------------
\newcommand{\R}{\mathbb{R}}
\newcommand{\bu}{\boldsymbol{u}}
\newcommand{\bv}{\boldsymbol{v}}
\newcommand{\bw}{\boldsymbol{w}}
\newcommand{\ba}{\boldsymbol{a}}
\newcommand{\bff}{\boldsymbol{f}}
\newcommand{\bze}{\boldsymbol{0}}
\newcommand{\ViscProj}{\Pi^{\mathrm{DS}}}
\newcommand{\norm}[1]{\lVert #1 \rVert}
\newcommand{\normK}[1]{\lVert #1 \rVert_{K}}
\newcommand{\tnorm}[1]{\lvert\!\lvert\!\lvert\, #1 \,\rvert\!\rvert\!\rvert}
\newcommand{\tns}{\tau_{1,\mathrm{NS}}}
\newcommand{\Eint}[2]{\mathrm{E}_{\mathrm{int},#1}(#2)}
\newcommand{\Dah}{\mathrm{Da}_{h,K}}
\newcommand{\Reh}{\mathrm{Re}_{h,K}}
\newcommand{\sigmatilde}{\widetilde{\sigma}_{\alpha}}
\newcommand{\Poneo}{\Pi_{1}^{\perp}}
\newcommand{\Ptwoo}{\Pi_{2}^{\perp}}
\newcommand{\smul}{\circledast}
\newcommand{\pmul}{\circ}

\newcommand{\forced}{\textsc{f}\xspace}   % forced by new physics
\newcommand{\tight}{\textsc{t}\xspace}    % tightening of exposition
\newcommand{\scope}{\textsc{s}\xspace}    % scope restriction

\title{Porosity-specific particularities of the OSGS a priori analysis\\[0.4ex]
\large A term-by-term comparison with Codina's orthogonal-subscale analysis of\\
the Oseen equations, and where the two proofs diverge}
\author{}
\date{\today}

\begin{document}
\maketitle

\begin{abstract}
This note compares, assumption by assumption and proof step by proof step, the
orthogonal-subgrid-scale (OSGS) a priori analysis of the linearized incompressible
inertial porous-media flow equations---the analysis carried out in the companion
manuscript~\cite{companion} and its self-contained backing note~\cite{osgsnote}---with
its direct ancestor, Codina's analysis of the Oseen equations by orthogonal
subscales~\cite{Codina2008}. Three questions are addressed. First, \emph{what is forced
by the new equation}: the porosity weight $\alpha(\boldsymbol x)$, the Darcy reaction
$\sigma\bu$, the artificial-compressibility term $\varepsilon p$, and the
deviatoric--symmetric viscous operator $\ViscProj$ each dictate concrete changes to the
hypotheses, the stabilization parameters, the working norm and the estimates, catalogued
in \cref{sec:forced}. Second, \emph{where the two proofs differ as proofs}, summarized as
a step-by-step ledger in \cref{sec:ledger}. Third---answering an explicit question, and
with the honest qualification that the answer is mostly ``no error, but the exposition can
be tightened''---\emph{where the Oseen analysis of~\cite{Codina2008} is less than fully
explicit}, and how the porous analysis sharpens it (\cref{sec:rigor}). Of the points
examined, exactly one (the smoothing lemma silently requiring a smooth-grading mesh
hypothesis) is a genuine, if benign, implicit assumption; the others are matters of
\emph{sharpness} (the $\ell^1$ error functional), \emph{disclosure} (stabilization
constants), or \emph{cleaner devices} (variable-$\tau$ projection handling). None is a
soundness defect: every theorem of~\cite{Codina2008} is correct as stated, and
R.~Codina is a co-author of~\cite{companion}. Throughout, ``Codina 2007'' denotes the
2007-online / 2008-print paper~\cite{Codina2008} (Appl.\ Numer.\ Math.\ \textbf{58}
(2008) 264--283); the companion note~\cite{osgsnote} cites it under the key
\texttt{Codina2008}.
\end{abstract}

\newpage
\tableofcontents
\newpage

% =====================================================================================
\section{Scope and the correct triangulation}
\label{sec:scope}

Codina 2007~\cite{Codina2008} analyzes the stationary \emph{Oseen} problem
\begin{equation}
  \label{eq:oseen}
  -\nu\Delta\bu + \ba\cdot\nabla\bu + \nabla p = \bff,
  \qquad \nabla\cdot\bu = 0 \ \text{ in }\Omega, \qquad \bu = \bze \ \text{ on }\partial\Omega,
\end{equation}
stabilized by orthogonal subscales. It contains \emph{three} analyzed formulations:
\textbf{Method~I} (a single projected residual per equation; Theorems~1--2), \textbf{Method~II}
(convective and pressure-gradient fluctuations controlled separately, in a finer norm;
Theorems~3--4), and a \textbf{pressure-only, viscous-dominated variant} whose norm depends
explicitly on the mesh and global Reynolds numbers (Theorems~5--6). The companion
paper~\cite{companion,osgsnote} treats the linearized porous problem
\begin{subequations}
\label{eq:porous}
\begin{alignat}{2}
  \alpha\,\ba\cdot\nabla\bu - 2\nabla\cdot(\alpha\nu\,\ViscProj\nabla\bu)
  + \alpha\nabla p + \sigma\bu &= \bff && \quad\text{in }\Omega, \\
  \varepsilon p + \nabla\cdot(\alpha\bu) &= 0 && \quad\text{in }\Omega,
  \qquad \bu=\bze \ \text{ on }\partial\Omega,
\end{alignat}
\end{subequations}
and adapts \textbf{only Method~I} (\cref{sec:ledger} records this as the single largest
scope difference). Method~II survives in~\cite{osgsnote} as a one-paragraph prospective
remark whose proof is not written out; the viscous-dominated pressure-only variant has no
porous analogue at all.

\paragraph{Two ancestor papers, two routes.}
A recurring source of confusion must be dispatched at the outset. The porous work has
\emph{two} a priori analyses following \emph{two different} Codina papers along
\emph{two different} routes:
\begin{center}
\begin{tabular}{@{}lll@{}}
\toprule
Porous variant & Template & Route \\
\midrule
ASGS (companion App.~C~\cite{companion}) & Codina 2001~\cite{Codina2001} & coercivity / C\'ea \\
OSGS (companion App.~D~\cite{companion}, $=$\cite{osgsnote}) & Codina 2007~\cite{Codina2008} & non-coercive inf--sup \\
\bottomrule
\end{tabular}
\end{center}
\emph{This note compares the OSGS row only: Codina 2007 versus the porous OSGS analysis.}
The consequence is important for a fair comparison of assumptions. Objects such as the
interior-face \emph{jump condition} on $\tau_1$, the \emph{enlarged inverse constant}
$\overline{C}_{\mathrm{inv}}$ for $\nabla\cdot(\alpha\nu\ViscProj\nabla\bu_h)$, and the
\emph{damped reaction coefficient} $\sigmatilde$ are \emph{Codina 2001} $\leftrightarrow$
\emph{porous-ASGS} objects. Codina 2007 is pure Oseen ($\sigma=0$) and has none of them; the
porous OSGS analysis \emph{dispenses} with the first two relative to the porous ASGS analysis
and never had the third. It would therefore be wrong to say ``the OSGS variant drops
hypothesis~H:jump of Codina 2007''---that hypothesis is not in Codina 2007.

\paragraph{A symbol collision to keep in view.}
Codina 2007 uses the letters $\alpha$ and $\sigma$ for \emph{design constants}:
$\alpha>1$ is the stabilization-amplification factor in $c_1=\alpha^2C_{\mathrm{inv}}^2$,
$c_2=\alpha C_{\mathrm{inv}}$, and $0<\sigma\le1$ appears in $c_3=\sigma$,
$c_4=\sigma/(\alpha C_{\mathrm{inv}})$ (his Eqs.~(27)--(28)). In the porous problem the
same letters are \emph{physical fields}: $\alpha$ is the porosity and $\sigma$ the Darcy
reaction. The porous note therefore renames Codina's amplification factor $\alpha\to\gamma$
(hypothesis (A6): $c_1\ge\gamma^2C_{\mathrm{inv}}^2$, $c_2\ge\gamma C_{\mathrm{inv}}$), and
Codina's design constant $\sigma$ simply disappears because, as \cref{sec:params} shows,
the porous $\tau_2$ is not free. \emph{One must not identify the porous $\alpha,\sigma$
with Codina's $\alpha,\sigma$.}

\paragraph{What is reused verbatim.}
The porous OSGS proof reuses from Codina 2007, essentially unchanged in \emph{structure}:
the non-coercive inf--sup argument built on a specially designed test function
($\tau\smash{\smul}\Pi_0$ of the projected residual); the continuous element-size surrogate
of his Lemma~1; the smoothing operation of his Lemma~2; the best-approximation property of
the weighted projection (his Remark~4, Eq.~(46)); and the ``interpolate--consistency--triangle''
skeleton of his Theorem~2. What changes is dictated, almost entirely, by the four new
ingredients of~\cref{eq:porous}.

% =====================================================================================
\section{The two model problems, side by side}
\label{sec:models}

\Cref{tab:models} isolates the four structural differences between~\cref{eq:oseen}
and~\cref{eq:porous}. Everything in \crefrange{sec:forced}{sec:ledger} is a consequence of
one of these four rows.

\begin{table}[htbp]
\centering
\small
\renewcommand{\arraystretch}{1.35}
\begin{tabular}{@{}p{0.20\linewidth}p{0.36\linewidth}p{0.36\linewidth}@{}}
\toprule
Ingredient & Codina 2007 (Oseen) & Porous linearized NS \\
\midrule
Viscous operator
 & $-\nu\Delta\bu$; energy $\nu\norm{\nabla\bu}^2$
 & $-2\nabla\cdot(\alpha\nu\ViscProj\nabla\bu)$; energy $2\nu\norm{\alpha^{1/2}\ViscProj\nabla\bu}^2$ \\
Coefficient weight
 & unweighted operators; $\nabla\cdot\ba=0$
 & everything weighted by porosity $\alpha(\boldsymbol x)$; weighted solenoidality $\nabla\cdot(\alpha\ba)=0$ \\
Zero-order term
 & none
 & Darcy reaction $\sigma\bu$, $\sigma\ge0$ constant \\
Mass equation
 & $\nabla\cdot\bu=0$ (exactly incompressible)
 & $\varepsilon p+\nabla\cdot(\alpha\bu)=0$ (artificial compressibility $\varepsilon\ge0$) \\
\bottomrule
\end{tabular}
\caption{The four ingredients that separate the porous problem from Oseen. The Galerkin
form gains, respectively, a deviatoric--symmetric viscous pairing, the porosity weight,
the reaction pairing $\sigma(\bv,\bu)$, and the compressibility pairing
$\varepsilon(q,p)$.}
\label{tab:models}
\end{table}

The Galerkin bilinear form and its coercivity (``diagonal'') identity make the four
additions concrete. In Codina 2007 (his Eqs.~(4), (39)),
\[
  B(U,V)=\nu(\nabla\bu,\nabla\bv)+(\ba\cdot\nabla\bu,\bv)-(p,\nabla\cdot\bv)+(q,\nabla\cdot\bu),
  \qquad
  B(U,U)=\nu\norm{\nabla\bu}^2;
\]
in the porous problem (\cite{osgsnote}, Eqs.~\texttt{GalerkinForm}, \texttt{CoercivityIdentity}),
\[
  B(\ba;U,V)=(\bv,\alpha\ba\cdot\nabla\bu)+2(\nabla\bv,\alpha\nu\ViscProj\nabla\bu)
  +(\bv,\alpha\nabla p)+\sigma(\bv,\bu)+\varepsilon(q,p)+(q,\nabla\cdot(\alpha\bu)),
\]
\[
  B(\ba;U,U)=2\nu\norm{\alpha^{1/2}\ViscProj\nabla\bu}^2+\norm{\sigma^{1/2}\bu}^2+\varepsilon\norm{p}^2 .
\]
The two \emph{new} diagonal quantities $\norm{\sigma^{1/2}\bu}^2$ and
$\varepsilon\norm{p}^2$, and the replacement of $\nu\norm{\nabla\bu}^2$ by the
deviatoric--symmetric $2\nu\norm{\alpha^{1/2}\ViscProj\nabla\bu}^2$, propagate into every
downstream estimate.

% =====================================================================================
\section{Assumptions: H1--H6 versus (A1)--(A8)}
\label{sec:assumptions}

\Cref{tab:assumptions} aligns Codina's six hypotheses with the porous eight and classifies
each porous item as \forced{} (forced by the new physics), \tight{} (a tightening or
explicit disclosure of something implicit in Codina 2007), or an inherited transcription.

\begin{table}[htbp]
\centering
\small
\renewcommand{\arraystretch}{1.3}
\begin{tabular}{@{}p{0.09\linewidth}p{0.40\linewidth}p{0.42\linewidth}@{}}
\toprule
Porous & Content & Relation to Codina 2007 \\
\midrule
(A1) & Nondegenerate mesh; inverse estimate; quasi-uniform patches; continuous element-size surrogate.
 & $=$ H3 $+$ his Lemma~1. Inherited. \\
(A2) & $\nu>0$, $\sigma\ge0$, $\varepsilon\ge0$ constant; Darcy tensor \emph{isotropic}, $\boldsymbol\sigma=\sigma\boldsymbol1$.
 & \forced. No Codina analogue; restricts the paper's general Forchheimer tensor $\boldsymbol\sigma(\alpha,\bu)$ so that the annihilation lemma holds. \\
(A3) & Porosity $\alpha\in W^{1,\infty}$, $0<\alpha_0\le\alpha\le\alpha_\infty\ (\le1)$; \emph{resolution} $h_K\norm{\nabla\alpha}_{\infty,K}\le C_{\nabla\alpha}\alpha_K$; and $\alpha\in W^{m,\infty}$, $m=\max(k_u,k_p)$, when $m\ge2$.
 & \forced. Entirely new. Codina has $\alpha\equiv1$. \\
(A4) & $\ba\in C^0$, elementwise $W^{1,\infty}$; \emph{weighted} solenoidality $\nabla\cdot(\alpha\ba)=0$; and $\norm{D^j\ba}_{\infty,K}\le C_D|\ba|_{\infty,K}$ for consistency.
 & \forced (weight) $+$ $=$ H1 (regularity) $+$ $=$ H2 (derivative bound). Codina: $\nabla\cdot\ba=0$. \\
(A5) & Continuous $V_h,Q_h$ of degrees $k_u,k_p\ge1$, $k_p\le k_u+1$; constants in $Q_h$; Dirichlet data.
 & Refines Codina's ``degree $k$ for both'': the porous analysis carries \emph{separate} $k_u,k_p$. \tight. \\
(A6) & $c_1\ge\gamma^2C_{\mathrm{inv}}^2$, $c_2\ge\gamma C_{\mathrm{inv}}$, $\gamma\ge\gamma_0$; and $\varepsilon\tau_{2,K}\le C_2\le1$.
 & $=$ H5 with amplification $\alpha\to\gamma$; \emph{explicit} admissible $\gamma_0$ (Codina leaves ``$\alpha$ large enough'' qualitative). $\varepsilon$-bound is new (\forced by $\varepsilon$). \tight/\forced. \\
(A7) & Projection stability: $\norm{z_h}_{\tau_1}\le\beta_0^{-1}(\norm{\Pi_{1,0}z_h}_{\tau_1}+\norm{\Poneo z_h}_{\tau_1})$ for $z_h=\alpha X(V_h)$.
 & $=$ H4 (his condition (35)), transported to the \emph{weighted} residual, and imposed only on the \emph{momentum} slot (see \cref{sec:ledger}). \\
(A8) & $\bu\in H^{k_u+1}\cap H^1_0$, $p\in H^{k_p+1}\cap Q_0$.
 & \tight. Stronger than H6, which assumes only $p\in H^{k}$ (one order below velocity). \\
\bottomrule
\end{tabular}
\caption{The porous hypotheses against Codina's. \forced{}: forced by the new physics.
\tight{}: a tightening/disclosure of something Codina left implicit or coarser.}
\label{tab:assumptions}
\end{table}

Two observations deserve emphasis.

\begin{observation}[The pressure-regularity/order refinement]
Codina's H6 asks only for $p\in H^{k}(\Omega)$ with a \emph{single} interpolation degree
$k$, and his error function carries the pressure at order $h^{k}$---one order below the
velocity's $h^{k+1}$. The porous (A8) asks for $p\in H^{k_p+1}$ and carries the pressure at
its own \emph{full} order $h^{k_p+1}$, with a possibly distinct degree $k_p\le k_u+1$. This
is a genuine hypothesis and error-order difference, orthogonal to the four physical
ingredients, and it is worth stating explicitly because it changes what ``optimal'' means
for the pressure.
\end{observation}

\begin{observation}[The reaction is not wholly unprecedented in Codina's programme]
Codina 2007 itself notes (immediately after its Eq.~(16)) that
reference~\cite{Codina2001}---the very paper the porous \emph{ASGS} analysis
follows---already analyzed ``zero-order terms coming from Coriolis forces and permeability
effects'', i.e.\ an algebraic-subscale analysis \emph{with} a permeability zero-order term.
Thus the Darcy reaction is unprecedented specifically in the \emph{orthogonal}-subscale
setting of Codina 2007, not in Codina's stabilized-FEM programme as a whole. The novelty
below is precise: it is the interaction of the reaction with the \emph{orthogonal
projection} (the annihilation of the reactive residual) that is new.
\end{observation}

% =====================================================================================
\section{How the new equation forces the differences}
\label{sec:forced}

We take the four ingredients in turn, tracing each to the assumptions, parameters, norm,
and estimates it dictates.

\subsection{Porosity weight \texorpdfstring{$\alpha(\boldsymbol x)$}{alpha(x)}}
\label{sec:alpha}

The porosity does not merely rescale; because it is a \emph{variable field} it changes the
algebra of every step.
\begin{enumerate}[leftmargin=1.5em,itemsep=0.3ex]
\item \emph{The stabilized residual carries the weight.} The natural stabilized quantity is
$\alpha X(V)=\alpha(\ba\cdot\nabla\bv+\nabla q)$, not the bare $X(V)$ of Codina. Both
stabilization terms $S_1,S_2$ (\cite{osgsnote}, Eq.~\texttt{Sdefs}) project the
\emph{weighted} momentum residual $\alpha X(U)$ and the \emph{weighted} divergence
$\nabla\cdot(\alpha\bu)$.
\item \emph{Weighted solenoidality.} The convective skew-symmetry that makes the Galerkin
form diagonal now requires $\nabla\cdot(\alpha\ba)=0$ (assumption (A4)):
$(\bu,\alpha\ba\cdot\nabla\bu)=\tfrac12(\alpha\ba,\nabla|\bu|^2)=-\tfrac12(\nabla\cdot(\alpha\ba),|\bu|^2)=0$.
Codina's $\nabla\cdot\ba=0$ (H1) is the $\alpha\equiv1$ special case.
\item \emph{Porosity-resolution condition (A3).} A hypothesis with \emph{no} Codina analogue:
$h_K\norm{\nabla\alpha}_{\infty,K}\le C_{\nabla\alpha}\alpha_K$. It is consumed in three
places: (i) a dedicated \emph{multiplicative} patch-equivalence lemma for $\alpha$
(\cite{osgsnote}, \texttt{lem:patch}), whose constant $\delta_\alpha^2=(1+C_{\nabla\alpha})^2$
depends only on $C_{\nabla\alpha}$ and \emph{not} on the porosity contrast
$\alpha_\infty/\alpha_0$; (ii) the $\nabla\alpha$ term that appears whenever
$\nabla\cdot(\alpha\bw)=\alpha\nabla\cdot\bw+\nabla\alpha\cdot\bw$ is expanded (stability
term T7, interpolation term I7); (iii) the Leibniz step in the weighted best-approximation
lemma. Codina's Lemma~1 equilibrates only $\tau$ across a patch; it never equilibrates a
\emph{coefficient} like $\alpha$.
\item \emph{The $\nabla\alpha$ term inflates the stability threshold.} Because
$\nabla\cdot(\alpha\bw_h)$ carries the extra $\nabla\alpha\cdot\bw_h$ piece, term T7 gains
the factor $M_2=1+C_{\nabla\alpha}/C_{\mathrm{inv}}$, which propagates into the admissible
design constant $\gamma_0=10(1+M_2)$. In Codina, where $\alpha\equiv1$, the corresponding
term has only $\nabla\cdot\bw_h$ and the threshold is merely ``$\alpha>3/2$''.
\item \emph{Weighted inverse estimates.} Codina's three closing inequalities on his p.~273 gain
the weight $\alpha_K$ and become $\mathrm{(K1)}$--$\mathrm{(K3)}$ of~\cite{osgsnote}
(e.g.\ $C_{\mathrm{inv}}(\alpha_K\nu)^{1/2}/h_K\le\gamma^{-1}\tau_1^{-1/2}$), with the
amplification $\alpha\to\gamma$ already noted.
\item \emph{Extra porosity smoothness for $k\ge2$.} The best-approximation lemma must project the
\emph{product} $\alpha z$; a Leibniz bound $\norm{\alpha z}_{H^r(S_K)}\le
C\,C_{\alpha,r}\alpha_{\infty}\norm{z}_{H^r(S_K)}$ forces the constant $C_{\alpha,r}$ and
hence the $W^{m,\infty}$ tail of (A3). Codina projects $z$ itself.
\end{enumerate}

\subsection{Darcy reaction \texorpdfstring{$\sigma\bu$}{sigma u}: the structural heart}
\label{sec:reaction}

This is where the analysis genuinely departs from Codina 2007, through a single mechanism.

\paragraph{Annihilation.} For constant $\sigma$, $\sigma\bu_h$ is a finite element function,
so its orthogonal fluctuation vanishes identically: $\Poneo(\sigma\bu_h)=\bze$
(\cite{osgsnote}, \texttt{lem:annihilation}). The orthogonal projection therefore
\emph{deletes} the reactive residual. There is no reactive subscale to renormalize the
reaction coefficient.

\paragraph{$\sigma$ versus $\sigmatilde$ (the trade).} Contrast the ASGS variant, whose
residual-based stabilization retains the reactive fluctuation and produces the diagonal
$\sigma(1-\sigma\tau_1)\norm{\bu_h}^2=\sigmatilde\norm{\bu_h}^2$ with the \emph{damped}
\[
  \sigmatilde=\frac{\tns^{-1}\sigma}{\tns^{-1}+\sigma/\alpha_K}\le\sigma .
\]
Because OSGS annihilates the reactive residual, no such damping occurs: the \emph{full}
$\norm{\sigma^{1/2}\bu_h}^2$ survives in the diagonal and in the working norm. This is a
strictly \emph{stronger} stability estimate---the porous OSGS triple norm dominates the ASGS
triple norm term by term (\cite{osgsnote}, \texttt{lem:normcomparison}). The price is paid
on the other side of the ledger: the reactive and pressure \emph{interpolation} errors, no
longer absorbable through a reactive subscale, must be bounded head-on, and the best
conversions available cost a Damk\"ohler factor.

\paragraph{The Damk\"ohler factor and its two new identities.} The reaction enters $\tau_1$
through its denominator, $\tau_1=(\alpha_K\tns^{-1}+\sigma)^{-1}$ (Codina's $\tau_1$ has no
reaction), and this forces two parameter inequalities with no Codina analogue, in terms of
the elemental \emph{mesh Damk\"ohler number} $\Dah=\sigma h_K^2/(\alpha_K\nu)$:
\begin{align}
  \text{(P5)}\quad & \alpha_K^2\tau_1\tau_2\ge\frac{h_K^2}{c_1+\Dah},
  \qquad\text{hence}\qquad \tau_2^{-1/2}\le(c_1+\Dah)^{1/2}\tfrac{\alpha_K}{h_K}\tau_1^{1/2}, \\
  \text{(P6)}\quad & \sigma^{1/2}\le\Dah^{1/2}\,\tfrac{\alpha_K}{h_K}\tau_2^{1/2}.
\end{align}
(P5) is the sole origin of the $(1+\Dah)^{1/2}$ inflation. When $\sigma=0$ it collapses to
Codina's exact relation ``$\tau_{1,K}$ behaves as $h_K^2\tau_{2,K}^{-1}$'' (the last line of
his Lemma~3 proof), i.e.\ $\alpha_K^2\tau_1\tau_2=h_K^2/c_1$. The reaction thus \emph{adds}
$\Dah$ to a quantity that was a clean equality in the Oseen case.

\paragraph{New proof terms.} The reaction produces the new stability pairing
$\mathrm{T3}=\sigma(\bu_h,\bw_h)$, absorbed by the full diagonal $\norm{\sigma^{1/2}\bu_h}^2$
using only $\sigma\tau_1\le1$; the new interpolation group I3; the reaction-inflated error
functional (\cref{sec:params}); and an entirely new corollary---the \emph{$\sigma$-robust
optimal $L^2$ velocity estimate} (\cite{osgsnote}, \texttt{cor:Ltwo}), obtained by dividing
the $\sigma$-controlled triple norm by $\sigma^{1/2}$, with the pressure contribution
\emph{decaying} as $(\sigma\nu)^{-1/2}$. Codina has no reaction and therefore no $L^2$
velocity control except in his Reynolds-degrading viscous-dominated variant.

\subsection{Artificial compressibility \texorpdfstring{$\varepsilon p$}{eps p}: a new consistency term}
\label{sec:eps}

The compressibility is easy to overlook but produces a qualitatively new phenomenon.
\begin{enumerate}[leftmargin=1.5em,itemsep=0.3ex]
\item \emph{New diagonal and new pairing.} The form gains $\varepsilon(q,p)$; the diagonal
gains $\varepsilon\norm{p_h}^2$; the stability expansion gains
$\mathrm{T4}=\varepsilon(p_h,r_h)$, controlled by the new design bound $\varepsilon\tau_2\le
C_2\le1$ (assumption (A6))---an upper bound on $\varepsilon$ absent from Codina.
\item \emph{The mass consistency term stops vanishing.} This is the salient point. In
Codina, $\nabla\cdot\bu=0$ \emph{exactly}, so the mass part of his consistency error (his
Eq.~(44), second term) is $(\Ptwoo(\nabla\cdot\bu),\cdot)=0$: it vanishes identically. In
the porous problem the exact mass equation reads $\nabla\cdot(\alpha\bu)=-\varepsilon p$, so
\[
  S_2(U,V_h)=-\varepsilon\bigl(\Ptwoo p,\ \nabla\cdot(\alpha\bv_h)\bigr)_{\tau_2}\neq0
  \quad(\varepsilon>0),
\]
a genuinely new consistency contribution that must be bounded head-on, via
$\varepsilon\le C_2\tau_2^{-1}$ and (P5), and contributes a $(c_1+\Dah)$-carrying term to the
error functional. It vanishes precisely when $\varepsilon=0$, isolating $\varepsilon$ as its
sole driver. There is also a matching interpolation group I4 and the annihilation
$\Ptwoo(\varepsilon p_h)=0$ for constant $\varepsilon$.
\end{enumerate}

\subsection{Deviatoric--symmetric viscous operator \texorpdfstring{$\ViscProj$}{Pi-DS}}
\label{sec:visc}

The viscous term is $-2\nabla\cdot(\alpha\nu\ViscProj\nabla\bu)$, with $\ViscProj$ the
pointwise orthogonal projection onto the deviatoric--symmetric part of a tensor. Two
consequences.
\begin{enumerate}[leftmargin=1.5em,itemsep=0.3ex]
\item \emph{Energy in a projected seminorm.} The diagonal and the norm carry
$2\nu\norm{\alpha^{1/2}\ViscProj\nabla\bu}^2$ instead of Codina's $\nu\norm{\nabla\bu}^2$.
The perturbation bound for the viscous cross term (T1, I1) uses $|\ViscProj\boldsymbol M|\le|\boldsymbol M|$
pointwise.
\item \emph{Definiteness of the norm is no longer trivial.} Because control is on
$\ViscProj\nabla\bv$, not $\nabla\bv$, the triple norm is definite only through a
Korn-type identity
\[
  \norm{\ViscProj\nabla\bv}^2=\tfrac12\norm{\nabla\bv}^2+\bigl(\tfrac12-\tfrac1d\bigr)\norm{\nabla\cdot\bv}^2,
\]
valid on $H^1_0$ (full Dirichlet data, via integration by parts), whence
$\norm{\nabla\bv}\le\sqrt2\,\norm{\ViscProj\nabla\bv}$; absent full Dirichlet data or
pressure normalization the triple norm is only a \emph{seminorm} (\cite{osgsnote},
\texttt{lem:definiteness}). Codina's $\nu^{1/2}\norm{\nabla\bv}$ is a norm outright by
Poincar\'e. The companion paper flags that this deviatoric structure is precisely what
would make a Korn-type inequality nontrivial for a full existence theory.
\item \emph{The neglected viscous residual carries $\nabla\alpha$ even at $k_u=1$.} In Codina,
the dropped viscous residual $\nu\Delta\bu_h$ vanishes on element interiors for linear
elements. Here the dropped residual is $2\nu(\ViscProj\nabla\bu_h)\nabla\alpha$, nonzero
already at $k_u=1$ whenever $\alpha$ varies (though small under (A3)). This slightly widens
the analyzed-versus-implemented distance of the Method~I simplification.
\end{enumerate}

\subsection{A parameter re-parametrization worth naming}
\label{sec:params}

Codina 2007 uses \emph{four independent} design constants: $\tau_1=[c_1\nu/h^2+c_2|\ba|/h]^{-1}$
and $\tau_2=c_3\nu+c_4|\ba|h$ (his Eqs.~(11)--(12)), with $c_1,c_2$ in (27) and
$c_3,c_4$ in (28). The porous $\tau_2$ is \emph{not free}: it is slaved to $\tns$,
\[
  \tau_{2,K}=\frac{h_K^2}{c_1\alpha_K\tns}=\frac{\nu}{\alpha_K}\Bigl(1+\tfrac{c_2}{c_1}\Reh\Bigr),
\]
so the exact identity $h_K^2\tns^{-1}=c_1\alpha_K\tau_2$ (P3) holds, and only $c_1,c_2$
survive as free constants. This coupling---not present in Codina---is what \emph{generates}
(P3) and (P5) and hence the Damk\"ohler factor. The reduction of the porous theorem to
Codina's Theorem~2 therefore requires the dictionary $c_3=1$, $c_4=c_2/c_1$ (in the limit
$\alpha\equiv1$, $\sigma=\varepsilon=0$); it is not a free relabeling of four constants but
a consequence of the two-constant structure.

The error functionals reflect the same story. Codina's is the $\ell^1$ element sum
(his Eq.~(36))
\[
  \varepsilon_{\mathrm C}(h)=\sum_K\Bigl(\tau_{1,K}^{-1/2}h_K^{k+1}\norm{\bu}_{k+1,K}
  +\tau_{2,K}^{-1/2}h_K^{k}\norm{p}_{k,K}\Bigr);
\]
the porous theorem is stated in the \emph{sharp} porosity-weighted broken $\ell^2$ functional
\[
  \Psi_{\mathrm O}(h)=\Bigl(\sum_K(c_1+\Dah)\tfrac{\alpha_K^2}{h_K^2}
  \bigl(\tau_{2,K}\Eint{K}{\bu}^2+\tau_{1,K}\Eint{K}{p}^2\bigr)\Bigr)^{1/2},
\]
whose coarser $\ell^1$ majorant $\mathcal E(h)$ (the ASGS error function times
$(1+\Dah)^{1/2}$) is retained only as a comparator. The relation $\Psi_{\mathrm O}\le
c_1^{1/2}\mathcal E$ and the deliberate choice to state the theorem in $\Psi_{\mathrm O}$ is
the subject of \cref{sec:rigor}\,(i).

% =====================================================================================
\section{Where the Oseen analysis is less than fully explicit}
\label{sec:rigor}

The question ``are there places where the original paper was not completely rigorous?'' has,
after a careful re-reading, a nuanced answer, and it is worth stating the verdict before the
evidence: \textbf{Codina 2007 contains no soundness error.} Every inequality is valid and
every theorem holds as stated. What the porous analysis does---appropriately, since Codina
co-authors it---is (i) state one estimate in a sharper norm, (ii) surface one implicit mesh
hypothesis, (iii) disclose the gap between provable and practical constants, and (iv) replace
one ``trivially'' with an exact, variable-coefficient-robust device. Only item (ii) is a
genuine (and entirely benign) implicit assumption. We take them in order, each with its fair
severity.

\subsection*{(i) The \texorpdfstring{$\ell^1$}{l1} error functional is coarser than the sharp \texorpdfstring{$\ell^2$}{l2} object (sharpness, not soundness)}

Codina's convergence Theorems~2 and~4 read $\tnorm{U-U_h}\le C\varepsilon_{\mathrm C}(h)$,
with $\varepsilon_{\mathrm C}(h)$ the \emph{$\ell^1$} element sum above. But the $\tau$-norm
$\norm{\cdot}_{\tau}=(\sum_K\tau_K\norm{\cdot}_K^2)^{1/2}$ that appears throughout the proof
is an \emph{$\ell^2$} object, and the passage from it to $\varepsilon_{\mathrm C}(h)$---his
Eq.~(46), $\norm{v-\Pi_\tau v}_\tau\le C\sum_K\tau_K^{1/2}h_K^r\norm{v}_{H^r(K)}$---uses the
valid but lossy step $(\sum b_K^2)^{1/2}\le\sum b_K$.

The gap is quantifiable. On a quasi-uniform mesh in $d$ dimensions with $N\sim h^{-d}$
elements, a smooth solution has $\norm{\bu}_{H^{k+1}(K)}=O(h^{d/2})$ (since $|K|=O(h^d)$).
With the diffusion-dominated scaling $\tau_1^{-1/2}\sim\nu^{1/2}/h$, the per-element term is
$a_K=O(\nu^{1/2}h^{k+d/2})$; the sharp $\ell^2$ functional is
$\Psi=(\sum a_K^2)^{1/2}=N^{1/2}a=O(\nu^{1/2}h^{k})$---the \emph{optimal} energy-norm
rate---while the literal $\ell^1$ sum is $\varepsilon_{\mathrm C}=\sum a_K=Na=O(\nu^{1/2}h^{k-d/2})$.
The ratio is exactly $N^{1/2}=h^{-d/2}$: for $d=2,k=1$, $\Psi=O(h)$ but
$\varepsilon_{\mathrm C}=O(1)$; for $d=3,k=1$, $\varepsilon_{\mathrm C}=O(h^{-1/2})$, which
does not even go to zero.

\begin{note}
This is a matter of \emph{sharpness}, not correctness. The theorem
$\tnorm{U-U_h}\le C\varepsilon_{\mathrm C}(h)$ is a true inequality; the optimal rate is
genuinely delivered by the $\ell^2$ object that already lives inside Codina's own proof
(the $\tau$-norm best-approximation). The $O(h^{k-d/2})$ reading is what one gets by
substituting local smooth-solution scalings term by term into the $\ell^1$ sum---something no
informed reader of the stabilized-FEM literature does, where $\sum_K h^{k+1}\norm{\bu}_{k+1,K}$
is understood at optimal order. The porous analysis simply makes this explicit: it states its
theorem in $\Psi_{\mathrm O}$ and \emph{names} the loss (``the $\ell^1$ passage to $\mathcal E$
can lose a mesh-cardinality factor $\sim h^{-d/2}$ on a quasi-uniform family, so the
optimal-order corollaries must be read off the $\ell^2$ functional''). Severity: negligible as
correctness, moderate as presentation. Note that the porous \emph{ASGS} appendix itself still
uses an $\ell^1$ error functional in the Codina style, so this refinement is specifically a
feature of the \emph{OSGS} note.
\end{note}

\subsection*{(ii) The smoothing lemma silently needs a smooth-grading mesh hypothesis (a genuine, benign gap)}

This is the one point at which a hypothesis is genuinely missing. Codina's Lemma~2 bounds
$\norm{\tau\pmul v_h-\tau\smul v_h}_K\le\tau_K\psi(h)\norm{v_h}_K$ with $\psi(h)\to0$, and
its proof reduces this to $\max_{a\in N_K}|\bar\tau_a/\tau_K-1|\to0$, then closes with the
single sentence: ``From the continuity assumed for the advection velocity $\ba$ we have that
$\bar\tau_a\to\tau_K$ as $h\to0$.'' But $\tau_1$ depends on $h$ through $c_1\nu/h^2+c_2|\ba|/h$,
with the element length the continuous surrogate $\chi_1(\boldsymbol x_K,h)$, while the nodal
$\bar\tau_a$ uses $\chi_1(\boldsymbol x_a,h)$. Writing $\rho=\chi_1(\boldsymbol x_a,h)/\chi_1(\boldsymbol x_K,h)$,
\[
  \frac{\bar\tau_{1,a}^{-1}}{\tau_{1,K}^{-1}}
  =\frac{c_1\nu/(\rho h_K)^2+c_2|\ba|/(\rho h_K)}{c_1\nu/h_K^2+c_2|\ba|/h_K}
  \ \xrightarrow{\ h\to0\ }\ \begin{cases}\rho^{-2}&\text{(viscous-dominated)}\\[0.3ex]\rho^{-1}&\text{(convection-dominated)}\end{cases}
\]
which equals $1$ \emph{only if} $\rho\to1$. Continuity of $\ba$ controls the amplitude
$|\ba(\boldsymbol x_a)|\to|\ba|_{\infty,K}$ but says nothing about $\rho$. His Lemma~1 delivers
only a \emph{bounded} ratio $\chi_2/\chi_1\le\chi_0<\infty$ (patch quasi-uniformity), and
bounded is not $\to1$. Indeed, even global quasi-uniformity with genuine size variation
(element sizes in $[\underline c\,h,\overline c\,h]$, $\overline c>\underline c$) can keep
$\rho$ away from $1$: a family with a fixed red/green refinement interface, refined uniformly,
is shape-regular and patch-quasi-uniform with $h$-independent constants, yet the neighbour
ratio stays $\approx2$ at the interface for every $h$, so in the viscous regime
$|\bar\tau_a/\tau_K-1|\approx|2^{-2}-1|=3/4\not\to0$. Thus Lemma~2 as written requires the
element-size function to have \emph{vanishing relative oscillation over each element}---a
smooth-grading / asymptotic local quasi-uniformity property strictly stronger than
nondegeneracy (H3), and not among H1--H6.

\begin{note}
The deduction is valid; only a hypothesis is left implicit, and the fix is one line (assume
smooth grading, or restrict to globally quasi-uniform families, on which $\chi_1$ is
essentially constant and the term is trivial---for a uniform mesh $\chi_1$ is constant and
$\rho\equiv1$). Codina's convergence theorems, invoked on such meshes, are untouched. The
porous smoothing lemma surfaces exactly this object: it decomposes $\Delta(\tau_1^{-1})$ into
contributions from $\Delta\alpha$, $\Delta|\ba|$, $\Delta(h^{-2})$, $\Delta(h^{-1})$ and
introduces a modulus $\omega_\chi(h)\to0$ for ``the relative oscillation of $\chi_1(\cdot,h)$
over an element, which vanishes as $h\to0$ under the smooth-grading assumption on the mesh
family (bounded patch equivalence alone does not force it)''. In full honesty, the porous note
flags $\omega_\chi(h)\to0$ inside the proof rather than promoting it to its standing
assumption list (A1); a maximally rigorous statement would make it a named hypothesis. This is
a rigor-bookkeeping tightening, not a correctness bug: it should credit Codina's construction,
not frame it as broken.
\end{note}

\subsection*{(iii) The stability constants are conditional and disconnected from those used in practice (disclosure)}

Codina's Theorem~1 holds ``for $h$ sufficiently small and $\alpha$ in (27)--(28) large
enough''. The proof exhibits the threshold only partially---it needs $\alpha>3/2$ for one
coercivity coefficient, plus positivity of several collected coefficients---and gives no
consolidated admissible $\alpha_0$. His numerical section then fixes $c_1=4$, $c_2=2$ with no
reconciliation to (27)--(28). These are self-consistent with $\alpha C_{\mathrm{inv}}=2$, i.e.\
$\alpha=2/C_{\mathrm{inv}}$; his own written step $\alpha>3/2$ would then require
$C_{\mathrm{inv}}<4/3$, which is not established and is generally false for higher-degree
elements (where $C_{\mathrm{inv}}$ grows). So the stability/convergence theorems are correctly
proved \emph{conditional existence} results (``there exists a large enough amplification''),
not validations of the $O(1)$ constants actually computed with.

\begin{note}
This is the standard, well-understood shape of stabilized/VMS a priori theory---provable
stability needs stabilization above a threshold typically far larger than, and disconnected
from, practice---and it is not a defect. The load-bearing point is one of \emph{disclosure}:
Codina 2007 does not flag the gap, whereas the porous OSGS analysis does, quantitatively. It
exhibits an explicit admissible $\gamma_0=10(1+M_2)\ (\ge20)$ and computes that this demands
$c_1\gtrsim400\,C_{\mathrm{inv}}^2$, $c_2\gtrsim20\,C_{\mathrm{inv}}$---\emph{well above} the
experimental $c_1=4k^4$ (or $16k^4$), $c_2=2k^2$, ``which are not shown to satisfy them.'' It
concludes that the result is ``a conditional inf--sup stability theorem, not a calibration or
validation of the stabilization constants employed in the computations.'' It likewise records that the
smallness threshold $h_0$ is not uniform in $\alpha_0$, $\nu$, or the moduli of continuity of
$\alpha$ and $\ba$. None of this contradicts robust convergence at $c_1=4$: the sufficient
conditions are explicitly non-sharp, so practice simply lies in the region the theory does not
certify. Severity as a criticism of Codina: low; value of the explicit disclosure: high.
\end{note}

\subsection*{(iv) ``Trivially \texorpdfstring{$\beta_0=1$}{beta0=1}'' and the weighted projection are exact only for globally constant \texorpdfstring{$\tau$}{tau} (cleaner devices)}

Two closely related points, both about \emph{variable} $\tau$, both handled more transparently
by the porous analysis, neither an error.

\emph{(a) The divergence slot.} Codina states that condition (35) ``is trivially verified with
$\beta_0=1$'' for $z_h=\nabla\cdot\bv_h$, because $\nabla\cdot\bv_h$ has zero mean and
therefore $\Pi_{\tau_2,0}(\nabla\cdot\bv_h)=\Pi_{\tau_2}(\nabla\cdot\bv_h)$. Testing the
weighted projection against the constant shows only that the \emph{$\tau_2$-weighted} integral
is preserved, not the \emph{unweighted} one; so for \emph{variable} $\tau_2$ with a continuous
pressure space the constrained and unconstrained projections differ by one mean-mode and the
sharp constant is $\beta_0<1$, not $1$. (Exactness holds when $\tau_2$ is \emph{globally}
constant, or when the pressure space is elementwise/discontinuous---not merely when $\tau_2$ is
elementwise-constant, which is already Codina's case.) The consequence is benign: (35)/H4 still
holds with a good $\beta_0\in(0,1]$, so his theorems are unaffected. The porous analysis reaches
``constant exactly one'' by a manifestly exact route instead: it projects
$\nabla\cdot(\alpha\bu_h)$ onto the \emph{unconstrained} $Q_h$ (an exact Pythagorean
decomposition, no middle mode) and restores admissibility with a mean shift $m_h$ of the
pressure \emph{test} function, shown invisible in every pairing (including the porous
$\varepsilon(p_h,r_h)$ term). This is why the porous note imposes the projection-stability
condition (A7) only on the \emph{momentum} slot and needs no analogue for the divergence.

\emph{(b) The weighted projection.} Codina's a priori proof requires the \emph{$\tau$-weighted}
projection (Remark~4: the weighted product is ``essential'' to the best-approximation
property~(46)), yet he presents the discrete matrix form only for globally constant $\tau_1$,
where $\Pi_{\tau_1}$ collapses to the plain $L^2$ projection ($M$ the ordinary Gram matrix).
For variable $\tau$ the implemented (plain-$L^2$) and analyzed (weighted) methods differ, and
the paper leaves this analyzed-versus-implemented distance implicit. The porous note makes it
explicit (its ``analyzed versus implemented'' remark, item~(i)): ``the analysis uses the
$\tau$-weighted projections; the implementation uses the plain $L^2$ projection; the two
coincide when $\tau_i$ is \emph{globally} constant.''

\begin{note}
Both are cleaner or more explicit treatments, not corrections: Codina's analyzed bilinear
form openly uses $\Pi_\tau$, his proof is correct for the method it analyzes, and his
constant-$\tau$ matrix form is transparently flagged. The porous mean-shift device would
equally sharpen the Oseen proof; $\varepsilon>0$ and the $\alpha$-weight only add cosmetic
terms ($\int\nabla\cdot(\alpha\bu_h)=0$ still). Present these as variable-coefficient-robust
reformulations that credit Codina's Remark~1, not as audit findings.
\end{note}

\subsection*{A point that is \emph{not} a gap: consistency framing}

For completeness, one candidate concern turned out to be an overclaim and should be recorded
as \emph{agreement}. Codina's phrase that dropping the viscous/force fluctuations ``leads to a
method which is still consistent (in a sense explained later; cf.\ Remark~3)'' is sometimes read
as loose. It is not. His Eq.~(44) states outright that ``the method is not consistent in the
classical variational sense, since the RHS of (44) is not zero''; his Lemma~3 bounds the
consistency error by $O(\varepsilon_{\mathrm C}(h))$; and his Remark~3 exhibits the fully
consistent \emph{extended} form $B^\star_I$ with auxiliary unknowns to which ``in a sense''
refers. The porous note (\texttt{lem:consistency}) \emph{transcribes} this---``not consistent
in the classical variational sense (cf.\ Codina Eq.~(44)); its consistency error is the value
of the stabilization terms at the exact solution''---rather than sharpening it. (One small
technical clarification: the Eq.-(44) non-consistency holds for \emph{all} $k\ge1$; the $k\ge2$
threshold governs only the separately-neglected viscous-projection term, which is zero for
$P_1$ because $\Delta\bu_h=0$.) The two analyses treat OSGS consistency identically and
rigorously.

% =====================================================================================
\section{Proof-step ledger: where the two OSGS proofs diverge}
\label{sec:ledger}

\Cref{tab:ledger} walks the two proofs in parallel and is the requested summary of ``places
in the OSGS proofs where there are differences worth mentioning.'' Read it as: \emph{same
skeleton, four physical perturbations, one scope restriction, and a handful of expository
tightenings}.

\begin{small}
\renewcommand{\arraystretch}{1.25}
\begin{longtable}{@{}p{0.16\linewidth}p{0.31\linewidth}p{0.44\linewidth}@{}}
\caption{Step-by-step ledger of the two OSGS proofs. \forced{}: forced by the new physics;
\tight{}: an expository tightening of Codina 2007; \scope{}: a scope restriction.}
\label{tab:ledger}\\
\toprule
Step & Codina 2007 & Porous OSGS (difference) \\
\midrule
\endfirsthead
\multicolumn{3}{@{}l}{\footnotesize\itshape \cref{tab:ledger} (continued)}\\
\toprule
Step & Codina 2007 & Porous OSGS (difference) \\
\midrule
\endhead
\bottomrule
\endfoot
Route &
Non-coercive inf--sup with a special test function (Thm~1). &
Same route; well-posedness obtained \emph{from} the inf--sup estimate on the square system.
\forced{} nothing---this is the inherited backbone. \\

Diagonal identity &
$B_I(U_h,U_h)=\nu\norm{\nabla\bu_h}^2+\norm{\Poneo\xi_h}_{\tau_1}^2+\norm{\Ptwoo\delta_h}_{\tau_2}^2$. &
Adds $\norm{\sigma^{1/2}\bu_h}^2$ and $\varepsilon\norm{p_h}^2$; viscous term becomes
$2\nu\norm{\alpha^{1/2}\ViscProj\nabla\bu_h}^2$; requires $\nabla\cdot(\alpha\ba)=0$. \forced. \\

$S_1,S_2$ &
Project $\ba\cdot\nabla\bu_h+\nabla p_h$ and $\nabla\cdot\bu_h$. &
Project the \emph{weighted} $\alpha X(U)$ and $\nabla\cdot(\alpha\bu)$; weighted projections
$\Pi_1,\Pi_2$. \forced. \\

Annihilation &
Only viscous/force fluctuations neglected (zero for $P_1$). &
Also: $\Poneo(\sigma\bu_h)=\bze$, $\Ptwoo(\varepsilon p_h)=0$ for constant $\sigma,\varepsilon$
(\texttt{lem:annihilation})---the reaction/compressibility residuals are deleted \emph{exactly}.
\forced; the mechanism of the whole note. \\

$\tau_1,\tau_2$ &
Four free constants $c_1,\dots,c_4$; $\tau_2=c_3\nu+c_4|\ba|h$. &
Reaction in the $\tau_1$ denominator; $\tau_2$ \emph{slaved} to $\tns$ (only $c_1,c_2$ free);
identity (P3); dictionary $c_3=1,c_4=c_2/c_1$ on reduction. \forced. \\

Smoothing &
Lemma~2, $\psi(h)\to0$ by continuity of $\ba$. &
Aggregates moduli of $\alpha$, $\ba$ \emph{and} the mesh-grading $\omega_\chi$; uses
$\alpha_0>0$; surfaces the smooth-grading requirement (\cref{sec:rigor}\,(ii)). \forced{}$+$\tight. \\

Patch equivalence &
Lemma~1 / Eq.~(31): equilibrates $\tau$ only. &
Adds a \emph{multiplicative} porosity equivalence (constant independent of the contrast
$\alpha_\infty/\alpha_0$); notes a multiplicative $|\ba|_\infty$ equivalence \emph{fails} near
zeros of $\ba$. \forced. \\

Best-approximation &
Eq.~(46) for $z$; Remark~4. &
\texttt{lem:bestapprox} for the \emph{product} $\alpha z$; Leibniz constant $C_{\alpha,r}$,
$W^{m,\infty}$ porosity. \forced. \\

Interpolant &
A single interpolant satisfying~(30). &
\emph{Two} operators: boundary-preserving nodal (drives the $H^1_0$ global IBP in I2/I5) and
Scott--Zhang of $\alpha z$ (drives the weighted best-approximation). \forced. \\

Triple norm / definiteness &
$\nu^{1/2}\norm{\nabla\bv}+\dots$; definite by Poincar\'e. &
Adds $2\nu\norm{\alpha^{1/2}\ViscProj\nabla\bv}^2$, $\norm{\sigma^{1/2}\bv}^2$,
$\varepsilon\norm{q}^2$; definiteness needs the Korn/$\ViscProj$ identity and \emph{full}
Dirichlet data (\texttt{lem:definiteness}); dominates the ASGS norm term by term. \forced. \\

Special test function &
$[\tau_1\smul\Pi_{\tau_1,0}\xi_h,\ \tau_2\smul\Pi_{\tau_2,0}\delta_h]$; both slots
\emph{constrained}; $\beta_0=1$ ``trivially'' for the divergence. &
Momentum slot constrained (needs (A7)); \emph{pressure} slot uses the \emph{unconstrained}
projection $+$ mean shift $m_h$, invisible in every term, making the mass-slot constant
\emph{exactly one}---no A7-analogue for the divergence. \tight{} (\cref{sec:rigor}\,(iv-a)). \\

Perturbation terms &
Viscous $+$ two mixed terms; no reaction, no compressibility, no $\nabla\alpha$. &
New pairings $\mathrm{T3}=\sigma(\bu_h,\bw_h)$ (absorbed by full $D_\sigma$),
$\mathrm{T4}=\varepsilon(p_h,r_h)$ (absorbed by $D_\varepsilon$, needs $\varepsilon\tau_2\le C_2$);
all inverse steps carry $\alpha$; T7 carries the $\nabla\alpha$ piece (factor $M_2$). \forced. \\

Collection / constants &
``$\alpha>3/2$'' suffices; no consolidated $\alpha_0$. &
Explicit $\gamma_0=10(1+M_2)$, $\psi(h_0)\le\min\{1/8,\beta_0^2/16\}$; declared
\emph{conditional}, not a validation of $c_1=4k^4$ (\cref{sec:rigor}\,(iii)). \tight. \\

Consistency &
Eq.~(44); mass part $\propto\Ptwoo(\nabla\cdot\bu)=0$ \emph{vanishes}. &
Momentum part weighted, splits $\tau_1^{-1/2}$ (Da enters); \emph{new} mass part
$S_2(U,V_h)=-\varepsilon(\Ptwoo p,\nabla\cdot(\alpha\bv_h))_{\tau_2}\neq0$. \forced{} (by $\varepsilon$). \\

Interpolation groups &
Lemma~4: six groups. &
Seven (I1--I7): new reactive I3 (via P6), compressibility I4 (via P5), plus the Damk\"ohler
factor entering the convective I2, compressibility I4 and pressure-gradient I5; I7 carries
$\nabla\alpha$. \forced. \\

Error functional &
Single $\ell^1$ sum $\varepsilon_{\mathrm C}(h)$, Eq.~(36). &
Sharp $\ell^2$ $\Psi_{\mathrm O}(h)$ with local $(c_1+\Dah)$; $\ell^1$ majorant $\mathcal E(h)$
demoted, $h^{-d/2}$ loss flagged (\cref{sec:rigor}\,(i)). \tight{}$+$\forced{} (Da). \\

Theorem $+$ corollaries &
Thm~2 (Method~I); \emph{also} Thms~3--4 (Method~II) and 5--6 (viscous-dominated). &
Only Method~I proved. \emph{New} corollaries: ASGS-norm domination up to $(c_1+\Dah)^{1/2}$,
and the $\sigma$-robust optimal $L^2$ velocity estimate. Method~II prospective only;
viscous-dominated variant absent. \forced{} (corollaries) $+$ \scope{} (Method~II, Thms~5--6). \\

Reduction check &
Thm~2 with $c_1{=}\alpha^2C_{\mathrm{inv}}^2$, etc. &
$\alpha\equiv1,\sigma{=}\varepsilon{=}0\Rightarrow\Dah{=}0$; $\Psi_{\mathrm O}\to c_1^{1/2}\varepsilon_{\mathrm C}$;
recovers Codina Thm~2 under the dictionary $c_3=1,c_4=c_2/c_1$. Consistency check. \\
\end{longtable}
\end{small}

% =====================================================================================
\section{Closing taxonomy}
\label{sec:taxonomy}

Every difference between the two OSGS analyses falls into exactly one of three classes.
\Cref{tab:taxonomy} is the one-page summary.

\begin{table}[htbp]
\centering
\small
\renewcommand{\arraystretch}{1.3}
\begin{tabular}{@{}p{0.24\linewidth}p{0.70\linewidth}@{}}
\toprule
Class & Differences \\
\midrule
\forced{} \ Forced by the new physics
 & Weighted residual $\alpha X$; weighted solenoidality $\nabla\cdot(\alpha\ba)=0$;
   porosity-resolution (A3) and patch-equivalence lemma; $\nabla\alpha$ term (factor $M_2$);
   deviatoric--symmetric energy and Korn definiteness; reaction diagonal and pairing T3;
   annihilation of the reactive/compressibility residual; the Damk\"ohler identities (P5),
   (P6) and factor $(1+\Dah)^{1/2}$; $\tau_2$ slaved to $\tns$; compressibility pairing T4
   and the \emph{new non-vanishing} mass consistency term; interpolation groups I3, I4, I7;
   the isotropic-constant-coefficient restriction (A2); the $\sigma$-robust $L^2$ corollary. \\
\tight{} \ Tightening / disclosure of Codina 2007
 & The sharp $\ell^2$ error functional with the $h^{-d/2}$ warning
   (\cref{sec:rigor}\,(i)); surfacing the smooth-grading modulus $\omega_\chi$
   (\cref{sec:rigor}\,(ii)); explicit conditional constants $\gamma_0$ and the
   theory-vs-practice disclosure (\cref{sec:rigor}\,(iii)); the exact
   variable-$\tau$ mass-slot device and the explicit weighted-vs-plain projection distinction
   (\cref{sec:rigor}\,(iv)); separate $k_u,k_p$ and the stronger pressure regularity (A8). \\
\scope{} \ Scope restriction
 & Only Method~I is transported. Method~II is a prospective remark (proof not written);
   the pressure-only viscous-dominated variant (Codina Thms~5--6) has no porous analogue;
   the coefficients are frozen constant/isotropic and the analysis is linear (given
   $\ba$, converged projection). \\
\bottomrule
\end{tabular}
\caption{Taxonomy of every catalogued difference. The overwhelming majority are \forced{}
by the four physical ingredients of~\cref{eq:porous}; a compact set of \tight{} items sharpen
Codina 2007's exposition (only the smooth-grading one being a genuine, benign implicit
hypothesis); the \scope{} items record what the porous analysis deliberately does not
attempt.}
\label{tab:taxonomy}
\end{table}

\paragraph{One-sentence conclusion.} The porous OSGS a priori analysis is a faithful,
strictly generalizing transcription of the Method~I analysis of Codina 2007, whose every
substantive deviation is forced by the porosity weight, the Darcy reaction, the artificial
compressibility, or the deviatoric--symmetric viscous operator; where it appears to
``correct'' the ancestor it is in fact only sharpening exposition or disclosing a standard
implicit convention, with the lone genuine---and entirely benign---finding being that the
smoothing lemma of both papers tacitly assumes a smoothly graded mesh family.

% =====================================================================================
\begin{thebibliography}{9}

\bibitem{companion}
G.~Casas, J.~Gonz\'alez-Us\'ua, R.~Codina, I.~de-Pouplana,
\emph{A stabilized finite element method for incompressible, inertial flows in
inhomogeneous porous media}, companion manuscript.

\bibitem{osgsnote}
\emph{Stability and a priori error analysis of the OSGS variant} (self-contained companion
note; identical to Appendix~D of~\cite{companion}), file
\texttt{theory/osgs\_a\_priori/osgs\_convergence.tex}.

\bibitem{Codina2008}
R.~Codina,
\emph{Analysis of a stabilized finite element approximation of the Oseen equations using
orthogonal subscales},
Appl.\ Numer.\ Math.\ \textbf{58} (2008) 264--283 (online 18 Jan.\ 2007).

\bibitem{Codina2001}
R.~Codina,
\emph{A stabilized finite element method for generalized stationary incompressible flows},
Comput.\ Methods Appl.\ Mech.\ Engrg.\ \textbf{190} (2001) 2681--2706.

\bibitem{CodinaBlasco1997}
R.~Codina, J.~Blasco,
\emph{A finite element formulation for the Stokes problem allowing equal velocity--pressure
interpolation},
Comput.\ Methods Appl.\ Mech.\ Engrg.\ \textbf{143} (1997) 373--391.

\bibitem{ScottZhang}
L.~R. Scott, S.~Zhang,
\emph{Finite element interpolation of nonsmooth functions satisfying boundary conditions},
Math.\ Comp.\ \textbf{54} (1990) 483--493.

\end{thebibliography}

\end{document}
